\documentclass[a4paper,11pt]{article}

\usepackage[french]{babel}
\usepackage[T1]{fontenc}
\usepackage[utf8]{inputenc}

\usepackage[a4paper,margin=2cm]{geometry}
\usepackage{amsmath,amssymb}
\usepackage{enumitem}
\usepackage{fancyhdr}
\usepackage{lastpage}

\pagestyle{fancy}
\fancyhf{}
\fancyfoot[C]{0501012610 --- Fiche conçue avec l'aide de ChatGPT (OpenAI).}
\renewcommand{\headrulewidth}{0pt}

\setlength{\parindent}{0pt}
\setlength{\parskip}{0.6em}

\begin{document}

{\large\textit{Réussir son entrée en 4\textsuperscript{e}}}

\vspace{0.3cm}

{\LARGE\bfseries Nombres décimaux et calculs}

\vspace{0.3cm}

\hrule

\vspace{0.6cm}

\textbf{Objectif}

Cette fiche a pour but de te redonner confiance dans les calculs sur les nombres décimaux. Tu apprendras à estimer un résultat avant de calculer, à choisir une méthode efficace et à résoudre des problèmes demandant un peu de réflexion.

\vspace{0.8cm}

\textbf{Exercice 1}

Contenu à venir...

\vspace{0.6cm}

\textbf{Exercice 2}

Contenu à venir...

\vspace{0.6cm}

\textbf{Exercice 3 $\star$}

Contenu à venir...

\vspace{0.6cm}

\textbf{Exercice 4 $\star\star$}

Contenu à venir...

\vfill

\hrule

\smallskip

\textit{Pour les curieux...}

Les nombres décimaux se sont développés progressivement en Europe à partir du XVI\textsuperscript{e} siècle. Ils ont permis de simplifier de nombreux calculs utilisés par les commerçants, les navigateurs et les scientifiques.

\end{document}

